% GENERATED FILE. Edit canonical lessons, then run npm run generate:books.
% canonical-source-hash: fnv1a64:ce8c423124b824a2
% canonical-lessons: TE-C462-kauntaru, TE-C462-sakha, TE-C462-purincu, TE-C462-jari-ceyu, TE-C462-manjuru-ceyu

\chapter{At the Office Counter}
\label{ch:te-at-the-office-counter}
\hlchaptermodality{462}

\begin{chapteropening}
\textbf{By the end of this chapter:} \emph{I can say a counter, a branch office, to fill in, to issue and to sanction.}

Complete the last lesson of chapter 462: I can say a counter, a branch office, to fill in, to issue and to sanction.
\end{chapteropening}

\section[\texorpdfstring{\te{కౌంటరు} (kauṇṭaru)}{కౌంటరు (kauṇṭaru)}]{\te{కౌంటరు} (kauṇṭaru) — a counter}

\label{lesson:TE-C462-kauntaru}



\begin{quote}
Before the new one: say the Telugu for to shift elsewhere, then the Telugu for to deliver.
\end{quote}

\subsection*{You'll want to know: \te{కౌంటరు}}
\textbf{\te{కౌంటరు}} — \emph{kauṇṭaru} — \textquotedblleft{}a counter\textquotedblright{}.

\begin{quote}
\te{టిక్కెట్టు} \te{కౌంటరు} \te{దగ్గర} \te{పెద్ద} \te{వరుస} \te{ఉంది}.
\end{quote}

\begin{quote}
\emph{ṭikkeṭṭu kauṇṭaru daggara pedda varusa undi.}
\end{quote}

\begin{quote}
There is a long queue at the ticket counter.
\end{quote}

Any service counter: in a bank, a station, a shop.

\begin{cousinweb}
\textbf{\te{కౌంటరు}} (\emph{kauṇṭaru}) is a loan from English \emph{counter}, with the Telugu final -\te{ఉ}.
\end{cousinweb}

\begin{grammarlens}[title={What you've built}]
One more word for getting around.
\end{grammarlens}

\subsection*{Your turn}
\begin{itemize}
\raggedright
  \item \emph{Say it:} \emph{kauṇṭaru}
  \item \emph{Say it:} \emph{kauṇṭaru}, once more
  \item \emph{Say it:} \emph{ṭikkeṭṭu kauṇṭaru daggara pedda varusa undi.} — There is a long queue at the ticket counter.
  \item \emph{Recall:} say the Telugu for to shift elsewhere, then the Telugu for to deliver, then say \emph{kauṇṭaru} again
\end{itemize}

\begin{tcolorbox}[breakable,colback=teal!4,colframe=teal!35!black,title={Before you move on}]
What does \te{కౌంటరు} mean? (\textquotedblleft{}A counter\textquotedblright{}.) Say it once more.
\end{tcolorbox}
\section[\texorpdfstring{\te{శాఖ} (śākha)}{శాఖ (śākha)}]{\te{శాఖ} (śākha) — a branch office; a branch}

\label{lesson:TE-C462-sakha}



\begin{quote}
Before the new one: say the Telugu for to deliver, then the Telugu for a counter.
\end{quote}

\subsection*{You'll want to know: \te{శాఖ}}
\textbf{\te{శాఖ}} — \emph{śākha} — \textquotedblleft{}a branch office; a branch\textquotedblright{}.

\begin{quote}
\te{ఆ} \te{బ్యాంకుకు} \te{మా} \te{పట్టణంలో} \te{కొత్త} \te{శాఖ} \te{వచ్చింది}.
\end{quote}

\begin{quote}
\emph{ā byāṅkuku mā paṭṭaṇaṁlō kotta śākha vaccindi.}
\end{quote}

\begin{quote}
That bank has opened a new branch in our town.
\end{quote}

A branch of a tree, of a bank, of a government: \te{వాతావరణ} \te{శాఖ}, the weather department.

\begin{cousinweb}
\textbf{\te{శాఖ}} (\emph{śākha}) is borrowed from Sanskrit \emph{śākhā}, \textquotedblleft{}a branch\textquotedblright{} of a tree (CDIAL 12376); Telugu uses it for the branches of offices too.
\end{cousinweb}

\begin{grammarlens}[title={What you've built}]
One more word for getting around.
\end{grammarlens}

\subsection*{Your turn}
\begin{itemize}
\raggedright
  \item \emph{Say it:} \emph{śākha}
  \item \emph{Say it:} \emph{śākha}, once more
  \item \emph{Say it:} \emph{ā byāṅkuku mā paṭṭaṇaṁlō kotta śākha vaccindi.} — That bank has opened a new branch in our town.
  \item \emph{Recall:} say the Telugu for to deliver, then the Telugu for a counter, then say \emph{śākha} again
\end{itemize}

\begin{tcolorbox}[breakable,colback=teal!4,colframe=teal!35!black,title={Before you move on}]
What does \te{శాఖ} mean? (\textquotedblleft{}A branch office; a branch\textquotedblright{}.) Say it once more.
\end{tcolorbox}
\section[\texorpdfstring{\te{పూరించు} (pūriñcu)}{పూరించు (pūriñcu)}]{\te{పూరించు} (pūriñcu) — to fill in; to complete}

\label{lesson:TE-C462-purincu}



\begin{quote}
Before the new one: say the Telugu for a counter, then the Telugu for a branch office.
\end{quote}

\subsection*{You'll want to know: \te{పూరించు}}
\textbf{\te{పూరించు}} — \emph{pūriñcu} — \textquotedblleft{}to fill in; to complete\textquotedblright{}.

\begin{quote}
\te{నేను} \te{దరఖాస్తు} \te{పత్రం} \te{పూరించి} \te{కౌంటరులో} \te{ఇచ్చాను}.
\end{quote}

\begin{quote}
\emph{nēnu darakhāstu patraṁ pūriñci kauṇṭarulō iccānu.}
\end{quote}

\begin{quote}
I filled in the application form and handed it in at the counter.
\end{quote}

Past (I …): \te{పూరించాను}, \emph{pūriñcānu}.

Said of a form, a blank or a gap: \te{ఖాళీలను} \te{పూరించు}, to fill the blanks.

\begin{cousinweb}
\textbf{\te{పూరించు}} (\emph{pūriñcu}) is built from Sanskrit \emph{pūr-}, \textquotedblleft{}to fill\textquotedblright{} (CDIAL 8335, \emph{pūrayati}, \textquotedblleft{}fills\textquotedblright{}), with the Telugu verb ending -\te{ఇంచు}.
\end{cousinweb}

\begin{grammarlens}[title={What you've built}]
One more word for getting around.
\end{grammarlens}

\subsection*{Your turn}
\begin{itemize}
\raggedright
  \item \emph{Say it:} \emph{pūriñcu}
  \item \emph{Say it:} \emph{pūriñcu}, once more
  \item \emph{Say it:} \emph{nēnu darakhāstu patraṁ pūriñci kauṇṭarulō iccānu.} — I filled in the application form and handed it in at the counter.
  \item \emph{Recall:} say the Telugu for a counter, then the Telugu for a branch office, then say \emph{pūriñcu} again
\end{itemize}

\begin{tcolorbox}[breakable,colback=teal!4,colframe=teal!35!black,title={Before you move on}]
What does \te{పూరించు} mean? (\textquotedblleft{}To fill in; to complete\textquotedblright{}.) Say it once more.
\end{tcolorbox}
\section[\texorpdfstring{\te{జారీ} \te{చేయు} (jārī cēyu)}{జారీ చేయు (jārī cēyu)}]{\te{జారీ} \te{చేయు} (jārī cēyu) — to issue}

\label{lesson:TE-C462-jari-ceyu}



\begin{quote}
Before the new one: say the Telugu for a branch office, then the Telugu for to fill in.
\end{quote}

\subsection*{You'll want to know: \te{జారీ} \te{చేయు}}
\textbf{\te{జారీ} \te{చేయు}} — \emph{jārī cēyu} — \textquotedblleft{}to issue\textquotedblright{}.

\begin{quote}
\te{ప్రభుత్వం} \te{కొత్త} \te{పాస్పోర్టు} \te{జారీ} \te{చేసింది}.
\end{quote}

\begin{quote}
\emph{prabhutvaṁ kotta pāspōrṭu jārī cēsindi.}
\end{quote}

\begin{quote}
The government issued a new passport.
\end{quote}

Past (I …): \te{జారీ} \te{చేశాను}, \emph{jārī cēśānu}.

Said of an office giving out a document, an order or a warning.

\begin{cousinweb}
\textbf{\te{జారీ} \te{చేయు}} (\emph{jārī cēyu}) is a Telugu phrase: \te{జారీ}, a loan from Arabic \emph{jārī}, \textquotedblleft{}running, current\textquotedblright{}, which came through Persian and Hindustani, and the native \te{చేయు}, \textquotedblleft{}to do\textquotedblright{} (DEDR 1957). To make a paper current is to issue it.
\end{cousinweb}

\begin{grammarlens}[title={What you've built}]
One more word for getting around.
\end{grammarlens}

\subsection*{Your turn}
\begin{itemize}
\raggedright
  \item \emph{Say it:} \emph{jārī cēyu}
  \item \emph{Say it:} \emph{jārī cēyu}, once more
  \item \emph{Say it:} \emph{prabhutvaṁ kotta pāspōrṭu jārī cēsindi.} — The government issued a new passport.
  \item \emph{Recall:} say the Telugu for a branch office, then the Telugu for to fill in, then say \emph{jārī cēyu} again
\end{itemize}

\begin{tcolorbox}[breakable,colback=teal!4,colframe=teal!35!black,title={Before you move on}]
What does \te{జారీ} \te{చేయు} mean? (\textquotedblleft{}To issue\textquotedblright{}.) Say it once more.
\end{tcolorbox}
\section[\texorpdfstring{\te{మంజూరు} \te{చేయు} (mañjūru cēyu)}{మంజూరు చేయు (mañjūru cēyu)}]{\te{మంజూరు} \te{చేయు} (mañjūru cēyu) — to sanction; to grant}

\label{lesson:TE-C462-manjuru-ceyu}



\begin{quote}
Before the new one: say the Telugu for to fill in, then the Telugu for to issue.
\end{quote}

\subsection*{You'll want to know: \te{మంజూరు} \te{చేయు}}
\textbf{\te{మంజూరు} \te{చేయు}} — \emph{mañjūru cēyu} — \textquotedblleft{}to sanction; to grant\textquotedblright{}.

\begin{quote}
\te{బ్యాంకు} \te{రైతులకు} \te{అప్పు} \te{మంజూరు} \te{చేసింది}.
\end{quote}

\begin{quote}
\emph{byāṅku raitulaku appu mañjūru cēsindi.}
\end{quote}

\begin{quote}
The bank granted a loan to the farmers.
\end{quote}

Past (I …): \te{మంజూరు} \te{చేశాను}, \emph{mañjūru cēśānu}.

An official word: a loan, a leave or a visa is \te{మంజూరు}, granted.

\begin{cousinweb}
\textbf{\te{మంజూరు} \te{చేయు}} (\emph{mañjūru cēyu}) is a Telugu phrase: \te{మంజూరు}, usually traced to Arabic \emph{manẓūr}, \textquotedblleft{}approved, accepted\textquotedblright{}, through Persian and Hindustani, and the native \te{చేయు}, \textquotedblleft{}to do\textquotedblright{} (DEDR 1957).
\end{cousinweb}

\begin{grammarlens}[title={What you've built}]
One more word for getting around.
\end{grammarlens}

\subsection*{Your turn}
\begin{itemize}
\raggedright
  \item \emph{Say it:} \emph{mañjūru cēyu}
  \item \emph{Say it:} \emph{mañjūru cēyu}, once more
  \item \emph{Say it:} \emph{byāṅku raitulaku appu mañjūru cēsindi.} — The bank granted a loan to the farmers.
  \item \emph{Recall:} say the Telugu for to fill in, then the Telugu for to issue, then say \emph{mañjūru cēyu} again
\end{itemize}

\begin{tcolorbox}[breakable,colback=teal!4,colframe=teal!35!black,title={Before you move on}]
What does \te{మంజూరు} \te{చేయు} mean? (\textquotedblleft{}To sanction; to grant\textquotedblright{}.) Say it once more.
\end{tcolorbox}
